% \iffalse meta-comment
%
% Copyright (C) 2019 by Zangwei Zheng <zhengzangw@gmail.com>
%
% This file may be distributed and/or modified under the conditions of
% the LaTeX Project Public License, either version 1.3c of this license
% or (at your option) any later version. The latest version of this 
% license is in:
%
%   http://www.latex-project.org/lppl.txt
%
% and version 1.3c or later is part of all distributions of LaTeX
% version 2005/12/01 or later.
%
% \fi
%
% \iffalse
%<*driver>
\ProvidesFile{njurepo.dtx}[2019/01/29 1.0.1 Nanjing University Report Template]
\documentclass{ltxdoc}
\usepackage{dtx-style}
    \EnableCrossrefs
    \CodelineIndex
    \RecordChanges
\begin{document}
    \DocInput{njurepo.dtx}
    \PrintChanges
    \PrintIndex
\end{document}
%</driver>
% \fi
%
% \CheckSum{0}
%
% \CharacterTable
%  {Upper-case    \A\B\C\D\E\F\G\H\I\J\K\L\M\N\O\P\Q\R\S\T\U\V\W\X\Y\Z
%   Lower-case    \a\b\c\d\e\f\g\h\i\j\k\l\m\n\o\p\q\r\s\t\u\v\w\x\y\z
%   Digits        \0\1\2\3\4\5\6\7\8\9
%   Exclamation   \!     Double quote  \"     Hash (number) \#
%   Dollar        \$     Percent       \%     Ampersand     \&
%   Acute accent  \'     Left paren    \(     Right paren   \)
%   Asterisk      \*     Plus          \+     Comma         \,
%   Minus         \-     Point         \.     Solidus       \/
%   Colon         \:     Semicolon     \;     Less than     \<
%   Equals        \=     Greater than  \>     Question mark \?
%   Commercial at \@     Left bracket  \[     Backslash     \\
%   Right bracket \]     Circumflex    \^     Underscore    \_
%   Grave accent  \`     Left brace    \{     Vertical bar  \|
%   Right brace   \}     Tilde         \~}
%
% \DoNotIndex{\newenvironment,\@bsphack,\@empty,\@esphack,\sfcode}
% \DoNotIndex{\addtocounter,\label,\let,\linewidth,\newcounter}
% \DoNotIndex{\noindent,\normalfont,\par,\parskip,\phantomsection}
% \DoNotIndex{\providecommand,\ProvidesPackage,\refstepcounter}
% \DoNotIndex{\RequirePackage,\setcounter,\setlength,\string,\strut}
% \DoNotIndex{\textbackslash,\texttt,\ttfamily,\usepackage}
% \DoNotIndex{\begin,\end,\begingroup,\endgroup,\par,\\}
% \DoNotIndex{\if,\ifx,\ifdim,\ifnum,\ifcase,\else,\or,\fi}
% \DoNotIndex{\let,\def,\xdef,\edef,\newcommand,\renewcommand}
% \DoNotIndex{\expandafter,\csname,\endcsname,\relax,\protect}
% \DoNotIndex{\Huge,\huge,\LARGE,\Large,\large,\normalsize}
% \DoNotIndex{\small,\footnotesize,\scriptsize,\tiny}
% \DoNotIndex{\normalfont,\bfseries,\slshape,\sffamily,\interlinepenalty}
% \DoNotIndex{\textbf,\textit,\textsf,\textsc}
% \DoNotIndex{\hfil,\par,\hskip,\vskip,\vspace,\quad}
% \DoNotIndex{\centering,\raggedright,\ref}
% \DoNotIndex{\c@secnumdepth,\@startsection,\@setfontsize}
% \DoNotIndex{\ ,\@plus,\@minus,\p@,\z@,\@m,\@M,\@ne,\m@ne}
% \DoNotIndex{\@@par,\DeclareOperation,\RequirePackage,\LoadClass}
% \DoNotIndex{\AtBeginDocument,\AtEndDocument}
%
% \changes{v1.0.0}{2019/01/22}{Initial version}
% \changes{v1.0.1}{2019/01/29}{Add more ability}
%
% \GetFileInfo{\jobname.dtx}
%
% \def\indexname{索引}
% \def\glossaryname{修改记录}
% \IndexPrologue{\section{\indexname}}
% \GlossaryPrologue{\section{\glossaryname}}

% \title{\bfseries\color{violet}\njurepo： 南京大学本科生范用报告}
% \author{郑奘巍 \\[5pt]\texttt{zhengzangw@gmail.com}}
% \date{\fileversion\ (\filedate)}
% \maketitle\thispagestyle{empty}
%
% \begin{abstract}\noindent
% 此宏包旨在建立一个免于配置的作业、实验报告模板。借鉴的宏包如下：
% \begin{itemize}
%  \item 清华大学\thuthesis https://github.com/xueruini/thuthesis
%  \item 南京大学 NJUBachelor https://github.com/ZLCao/NJUBachelor
% \end{itemize}
% \end{abstract}
%
%
% \vskip2cm
% \def\abstractname{免责声明}
% \begin{abstract}
% \noindent
% \begin{enumerate}
% \item 本模板的发布遵守 \LaTeX\ Project Public License，使用前请认真阅读协议内
%   容。
% \item 本模板为作者根据清华大学的论文模板\thuthesis{}等改写而成。因作者能力有限，可以选择使用\thuthesis{}模板
% 并自定义选项和手动添加宏包，可能达到比本模板更加好的效果。
% \item 任何个人或组织以本模板为基础进行修改、扩展而生成的新的专用模板，请严格遵
%   守 \LaTeX\ Project Public License 协议。由于违犯协议而引起的任何纠纷争端均与
%   本模板作者无关。
% \end{enumerate}
% \end{abstract}
%
% \clearpage
% \pagestyle{fancy}
% \begin{multicols}{2}[
%   \setlength{\columnseprule}{.4pt}
%   \setlength{\columnsep}{18pt}]
%   \tableofcontents
% \end{multicols}
% \clearpage
%
%
% \section{安装}
% \label{sec:installation}
% \njurepo 尚未包含于\TeX 发行版中，需要自行前往ginjub主页安装。相关链接:
% \begin{itemize}
% \item 主页: 有待添加
% \end{itemize}
%
% \subsection{模板的组成}
% 下表列出了\njurepo 的主要文件及其功能介绍：
%
% \begin{longtable}{l|p{8cm}}
% \toprule
% {\heiti 文件（夹）} & {\heiti 功能描述}\\\midrule
% \endfirsthead
% \midrule
% {\heiti 文件（夹）} & {\heiti 功能描述}\\\midrule
% \endhead
% \endfoot
% \endlastfoot
% njurepo.ins & \textsc{DocStrip} 驱动文件（开发用） \\
% njurepo.dtx & \textsc{DocStrip} 源文件（开发用）\\\midrule
% njurepo.cls & 模板类文件\\\midrule
% njurepo.sty & 为示例文档加载其它宏包\\\midrule
% \textbf{njurepo.pdf} & 用户手册（本文档）\\\bottomrule
% \end{longtable}
%
% \section{使用说明}
% \subsection{选项}
% \label{sec:option}
% \DescribeOption{language}
% 论文的主要语言（默认：中文）。可选：\option{chinese}，\option{english}。
% 
% \DescribeOption{secret}
% \DescribeOption{open}
% 尚未添加的选择
% \begin{itemize}
% \item secret: 同thuthesis的设置
% \item open: any为页面连续，right为章节奇数页
% \end{itemize}
%
% \subsection{字体配置}
% \label{sec:font-config}
% 使用\CTeX\ 默认字体配置
% \subsubsection{字体命令}
% \label{sec:fontcmds}
% \myentry{字体}
% \DescribeMacro{\songti}
% \DescribeMacro{\fangsong}
% \DescribeMacro{\heiti}
% \DescribeMacro{\kaishu}
% 用来切换宋体、仿宋、黑体、楷体四种基本字体。
% \myentry{字号}
% \DescribeMacro{\chuhao}
% \DescribeMacro{\xiaochu}
% \DescribeMacro{\yihao}
% \DescribeMacro{\xiaoyi}
% \DescribeMacro{\bahao}
% 定义字体大小，分别为
% \begin{center}
% \begin{tabular}{llllll}
% \toprule
% \cs{chuhao} & \cs{xiaochu} & \cs{yihao}  & \cs{xiaoyi}    & \cs{erhao}  & \cs{xiaoer}\\
% \cs{sanhao} & \cs{xiaosan} & \cs{sihao}  & \cs{banxiaosi} & \cs{xiaosi} & \cs{dawu}\\
% \cs{wuhao}  & \cs{xiaowu}  & \cs{liuhao} & \cs{xiaoliu}   & \cs{qihao}  & \cs{bahao}\\\bottomrule
% \end{tabular}
% \end{center}
% 使用方法为：\cs{command}\oarg{num}，其中 command 为字号命令，num 为行距。比
% 如 \cs{xiaosi}[1.5] 表示选择小四字体，行距 1.5 倍。写作指南要求表格中的字体
% 是 \cs{dawu}，模板已经设置好了。
%
% \StopEventually{}
%
% \section{实现细节}
% \subsection{基本信息}
%    \begin{macrocode}
%<*cls>
\NeedsTeXFormat{LaTeX2e}
\ProvidesClass{njurepo}[2019/01/25 1.0.0 Nanjing University Report Template]
%    \end{macrocode}
%
% \subsection{定义选项}
% \label{sec:defoption}
% 使用kvoptions宏包进行选项设置
%    \begin{macrocode}
\hyphenation{NJU-repo}
\def\njurepo{\textsc{NJU}\-\textsc{repo}}
\def\thuthesis{\textsc{Thu}\-\textsc{Thesis}}
\def\version{1.0.1}
\RequirePackage{kvoptions}
\SetupKeyvalOptions{
    family=nju,
    prefix=nju@,
    setkeys=\kvsetkeys
}
\DeclareStringOption[chinese]{language}[chinese]
\DeclareStringOption[any]{open}[any]
\DeclareBoolOption{wide}
\DeclareBoolOption{color}
\DeclareBoolOption{draft}
\DeclareBoolOption{awesomefont}
\DeclareDefaultOption{\PassOptionsToClass{\CurrentOption}{ctexbook}}

\ProcessKeyvalOptions*
%    \end{macrocode}
%
% 检测选项是否合法
%    \begin{macrocode}
\newcommand\nju@validate@key[1]{%
  \@ifundefined{nju@\csname nju@#1\endcsname true}{%
    \ClassError{njurepo}{Invalid value '\csname nju#1\endcsname'}{}
    }{%
      \csname nju@\csname nju@#1\endcsname true\endcsname
    }
}
\newif\ifnju@chinese
\newif\ifnju@english
\nju@validate@key{language}
\newif\ifnju@any
\newif\ifnju@right
\nju@validate@key{open}
%    \end{macrocode}
% 
% 使用ctexbook宏包
%    \begin{macrocode}
\LoadClass[a4paper,openany,UTF8,zihao=-4,scheme=plain]{ctexbook}
%    \end{macrocode}
%
% \subsection{加载宏包}
% \label{sec:loadpkg}
% 用于开发的宏包
%    \begin{macrocode}
\RequirePackage{etoolbox}
\RequirePackage{ifxetex}
\RequirePackage{xparse}
\RequirePackage{amsmath}
\RequirePackage{graphicx}
\graphicspath{{figs/}}
\graphicspath{{figs/logo/}}
\RequirePackage[labelformat=simple]{subcaption}
\RequirePackage{pdfpages}
\includepdfset{fitpaper=true}
\RequirePackage{CJKfntef}
\RequirePackage[amsmath, thmmarks, hyperref]{ntheorem}
\RequirePackage{array}
\RequirePackage{longtable}
\RequirePackage[sort&compress]{natbib}
\RequirePackage{booktabs}
%    \end{macrocode}
%
% 超链接插入
%    \begin{macrocode}
\RequirePackage{hyperref}
\ifxetex
  \hypersetup{%
    CJKbookmarks=true}
\else
  \hypersetup{%
    unicode=true,
    CJKbookmarks=false}
\fi
\hypersetup{%
  linktoc=all,
  bookmarksnumbered=true,
  bookmarksopen=true,
  bookmarksopenlevel=1,
  breaklinks=true,
  colorlinks=false,
  plainpages=false,
  pdfborder=0 0 0}	
\urlstyle{same}
\def\UrlBreaks{%
  \do\/%
  \do\a\do\b\do\c\do\d\do\e\do\f\do\g\do\h\do\i\do\j\do\k\do\l%
     \do\m\do\n\do\o\do\p\do\q\do\r\do\s\do\t\do\u\do\v\do\w\do\x\do\y\do\z%
  \do\A\do\B\do\C\do\D\do\E\do\F\do\G\do\H\do\I\do\J\do\K\do\L%
     \do\M\do\N\do\O\do\P\do\Q\do\R\do\S\do\T\do\U\do\V\do\W\do\X\do\Y\do\Z%
  \do0\do1\do2\do3\do4\do5\do6\do7\do8\do9\do=\do/\do.\do:%
  \do\*\do\-\do\~\do\'\do\"\do\-}
\Urlmuskip=0mu plus 0.1mu
%    \end{macrocode}
%
% 页眉页脚设置
%    \begin{macrocode}
\RequirePackage{fancyhdr}
\RequirePackage{notoccite}	
%    \end{macrocode}
%
% \subsection{页面设置}
% 使用了thuthesis的非本科生默认配置。
%    \begin{macrocode}
\RequirePackage{geometry}
\ifnju@wide 
\geometry{
    a4paper, %210*297mm
    hcentering,
    ignoreall,
    nomarginpar,
    left=10mm,
    headheight=5mm,
    headsep=5mm,
    textheight=237mm,
    bottom=29mm,
    footskip=6mm
}\else
\geometry{
    a4paper, %210*297mm
    hcentering,
    ignoreall,
    nomarginpar,
    left=30mm,
    headheight=5mm,
    headsep=5mm,
    textheight=237mm,
    bottom=29mm,
    footskip=6mm
}
\fi
%    \end{macrocode}
%
% \subsection{主文档格式}
% \label{sec:mainbody}
%
% \begin{macro}{\cleardoublepage}
%    \begin{macrocode}
\let\nju@cleardoublepage\cleardoublepage
\newcommand{\nju@clearemptydoublepage}{%
  \clearpage{\pagestyle{nju@empty}\nju@cleardoublepage}}
\let\cleardoublepage\nju@clearemptydoublepage
%    \end{macrocode}
% \end{macro}
%
% \begin{macro}{\frontmatter}
% \begin{macro}{\mainmatter}
% \begin{macro}{\backmatter}
%    \begin{macrocode}
\renewcommand\frontmatter{%
    \ifnju@right\cleardoublepage\else\clearpage\fi
    \@mainmatterfalse
    \pagenumbering{Roman}
    \pagestyle{nju@empty}}
\renewcommand\mainmatter{%
    \ifnju@right\cleardoublepage\else\clearpage\fi
    \@mainmattertrue
    \pagenumbering{arabic}
    \pagestyle{nju@headings}}
\renewcommand\backmatter{%
    \ifnju@right\cleardoublepage\else\clearpage\fi
    \@mainmattertrue}
%    \end{macrocode}
% \end{macro}
% \end{macro}
% \end{macro}
%
% \subsection{字体与字号}
% \label{sec:font}
% \subsubsection{英文字体}
% 配置英文字体。
%    \begin{macrocode}
\newcommand\nju@fontset{\csname g__ctex_fontset_tl\endcsname}
\ifthenelse{\equal{\nju@fontset}{fandol}}{
  \setmainfont[
    Extension      = .otf,
    UprightFont    = *-regular,
    BoldFont       = *-bold,
    ItalicFont     = *-italic,
    BoldItalicFont = *-bolditalic,
  ]{texgyretermes}
  \setsansfont[
    Extension      = .otf,
    UprightFont    = *-regular,
    BoldFont       = *-bold,
    ItalicFont     = *-italic,
    BoldItalicFont = *-bolditalic,
  ]{texgyreheros}
  \setmonofont[
    Extension      = .otf,
    UprightFont    = *-regular,
    BoldFont       = *-bold,
    ItalicFont     = *-italic,
    BoldItalicFont = *-bolditalic,
    Scale          = MatchLowercase,
  ]{texgyrecursor}
}{
  \setmainfont{Times New Roman}
  \setsansfont{Arial}[SmallCapsFont=spinweradCBold]
  \ifthenelse{\equal{\nju@fontset}{mac}}{
    \setmonofont[Scale=MatchLowercase]{Menlo}
  }{
    \setmonofont[Scale=MatchLowercase]{Courier New}
  }
}
%    \end{macrocode}
%
% \subsubsection{数学环境字体}
% 配置数学字体（使用unicode-math）
%    \begin{macrocode}
\RequirePackage{unicode-math}
\unimathsetup{
  math-style = ISO,
  bold-style = ISO,
  nabla      = upright,
  partial    = upright,
}
\IfFontExistsTF{STIX2Math.otf}{
  \setmathfont[StylisticSet=8]{STIX2Math.otf}
  \setmathfont[range={scr,bfscr},StylisticSet=1]{STIX2Math.otf}
  \IfFontExistsTF{XITSMath-Regular.otf}{
    \setmathfont[range={\partial,\lbrace,\rbrace}]{XITSMath-Regular.otf}
  }{
    \setmathfont[range={\partial,\lbrace,\rbrace}]{xits-math.otf}
  }
}{
  \setmathfont[
    Extension    = .otf,
    BoldFont     = *bold,
    StylisticSet = 8,
  ]{xits-math}
  \setmathfont[range={cal,bfcal},StylisticSet=1]{xits-math.otf}
}
%    \end{macrocode}
%
% \subsubsection{数学环境符号}
% \begin{macro}{\ldots}
% 省略号一律居中，所以 \cs{ldots} 不再居于底部。
%    \begin{macrocode}
\ifnju@chinese
  \def\mathellipsis{\cdots}
\fi
%    \end{macrocode}
% \end{macro}
%
% \begin{macro}{\le}
% \begin{macro}{\ge}
% \begin{macro}{\leq}
% \begin{macro}{\geq}
% 小于等于号要使用倾斜的形式。
%    \begin{macrocode}
\protected\def\le{\leqslant}
\protected\def\ge{\geqslant}
\AtBeginDocument{%
  \renewcommand\leq{\leqslant}%
  \renewcommand\geq{\geqslant}%
}
%    \end{macrocode}
% \end{macro}
% \end{macro}
% \end{macro}
% \end{macro}
%
% \begin{macro}{\int}
% 积分号 \cs{int} 使用正体，并且上下限默认置于积分号上下两侧。
%    \begin{macrocode}
\removenolimits{%
  \int\iint\iiint\iiiint\oint\oiint\oiiint
  \intclockwise\varointclockwise\ointctrclockwise\sumint
  \intbar\intBar\fint\cirfnint\awint\rppolint
  \scpolint\npolint\pointint\sqint\intlarhk\intx
  \intcap\intcup\upint\lowint
}
%    \end{macrocode}
% \end{macro}
%
% \begin{macro}{\Re}
% \begin{macro}{\Im}
% \begin{macro}{\nabla}
% 实部、虚部操作符使用罗马体 $\mathrm{Re}$、$\mathrm{Im}$ 而不是 fraktur 体
% $\Re$、$\Im$。\cs{nabla} 使用粗正体。
%    \begin{macrocode}
\AtBeginDocument{%
  \renewcommand{\Re}{\operatorname{Re}}%
  \renewcommand{\Im}{\operatorname{Im}}%
  \renewcommand\nabla{\mbfnabla}%
}
%    \end{macrocode}
% \end{macro}
% \end{macro}
% \end{macro}
%
% \begin{macro}{\bm}
% \begin{macro}{\boldsymbol}
% 兼容旧的粗体命令：\pkg{bm} 的 \cs{bm} 和 \pkg{amsmath} 的 \cs{boldsymbol}。
%    \begin{macrocode}
\newcommand\bm{\symbf}
\renewcommand\boldsymbol{\symbf}
%    \end{macrocode}
% \end{macro}
% \end{macro}
%
% \begin{macro}{\square}
% 兼容 \pkg{amssymb} 中的命令。
%    \begin{macrocode}
\newcommand\square{\mdlgwhtsquare}
%    \end{macrocode}
% \end{macro}
%
% 允许太长的公式断行、分页等。
%    \begin{macrocode}
\allowdisplaybreaks[4]
\renewcommand\theequation{\ifnum \c@chapter>\z@ \thechapter-\fi\@arabic\c@equation}
%    \end{macrocode}
%
% 公式距前后文的距离由 4 个参数控制，参见 \cs{normalsize} 的定义。
%    \begin{macrocode}
\def\make@df@tag{\@ifstar\nju@make@df@tag@@\make@df@tag@@@}
\def\nju@make@df@tag@@#1{\gdef\df@tag{\nju@maketag{#1}\def\@currentlabel{#1}}}
\def\nju@maketag#1{\maketag@@@{(\ignorespaces #1\unskip\@@italiccorr)}}
\def\tagform@#1{\maketag@@@{(\ignorespaces #1\unskip\@@italiccorr)\equcaption{#1}}}
%    \end{macrocode}
% 修改 \cs{tagform} 会影响 \cs{eqref}。
%    \begin{macrocode}
\renewcommand{\eqref}[1]{\textup{(\ref{#1})}}
%    \end{macrocode}
%
% \subsubsection{中文字体}
% \pkg{ctex}在微软下使用雅黑字体，在macOS下使用苹方字体。这里不做更改。
%
% \subsubsection{字号}
% WORD 中的字号对应该关系如下（1bp = 72.27/72 pt）:
% \begin{center}
% \begin{tabular}{llll}
% \toprule
% 初号 & 42bp & 14.82mm & 42.1575pt \\
% 小初 & 36bp & 12.70mm & 36.135 pt \\
% 一号 & 26bp & 9.17mm & 26.0975pt \\
% 小一 & 24bp & 8.47mm & 24.09pt \\
% 二号 & 22bp & 7.76mm & 22.0825pt \\
% 小二 & 18bp & 6.35mm & 18.0675pt \\
% 三号 & 16bp & 5.64mm & 16.06pt \\
% 小三 & 15bp & 5.29mm & 15.05625pt \\
% 四号 & 14bp & 4.94mm & 14.0525pt \\
% 小四 & 12bp & 4.23mm & 12.045pt \\
% 五号 & 10.5bp & 3.70mm & 10.59375pt \\
% 小五 & 9bp & 3.18mm & 9.03375pt \\
% 六号 & 7.5bp & 2.56mm & \\
% 小六 & 6.5bp & 2.29mm & \\
% 七号 & 5.5bp & 1.94mm & \\
% 八号 & 5bp & 1.76mm & \\\bottomrule
% \end{tabular}
% \end{center}
%
% \begin{macro}{\normalsize}
% 默认正文小四号 (12bp) 字，行距为固定值 20 bp。
%    \begin{macrocode}
\renewcommand\normalsize{%
  \@setfontsize\normalsize{12bp}{20bp}%
  \abovedisplayskip=12bp \@plus 2bp \@minus 2bp
  \abovedisplayshortskip=12bp \@plus 2bp \@minus 2bp
  \belowdisplayskip=\abovedisplayskip
  \belowdisplayshortskip=\abovedisplayshortskip}
%    \end{macrocode}
% \end{macro}
%
% \begin{macro}{\nju@def@fontsize}
% 根据习惯定义字号。用法：
% \cs{nju@def@fontsize}\marg{字号名称}\marg{磅数}
%
% 避免了字号选择和行距的紧耦合。所有字号定义时为单倍行距，并提供选项指定行距倍数。
%    \begin{macrocode}
\def\nju@def@fontsize#1#2{%
  \expandafter\newcommand\csname #1\endcsname[1][1.3]{%
    \fontsize{#2}{##1\dimexpr #2}\selectfont}}
%    \end{macrocode}
% \end{macro}
%
% \begin{macro}{\chuhao}
% \begin{macro}{\xiaochu}
% \begin{macro}{\yihao}
% \begin{macro}{\xiaoyi}
% \begin{macro}{\erhao}
% \begin{macro}{\xiaoer}
% \begin{macro}{\sanhao}
% \begin{macro}{\xiaosan}
% \begin{macro}{\sihao}
% \begin{macro}{\banxiaosi}
% \begin{macro}{\xiaosi}
% \begin{macro}{\dawu}
% \begin{macro}{\wuhao}
% \begin{macro}{\xiaowu}
% \begin{macro}{\liuhao}
% \begin{macro}{\xiaoliu}
% \begin{macro}{\qihao}
% \begin{macro}{\bahao}
% 一组字号定义。
%    \begin{macrocode}
\nju@def@fontsize{chuhao}{42bp}
\nju@def@fontsize{xiaochu}{36bp}
\nju@def@fontsize{yihao}{26bp}
\nju@def@fontsize{xiaoyi}{24bp}
\nju@def@fontsize{erhao}{22bp}
\nju@def@fontsize{xiaoer}{18bp}
\nju@def@fontsize{sanhao}{16bp}
\nju@def@fontsize{xiaosan}{15bp}
\nju@def@fontsize{sihao}{14bp}
\nju@def@fontsize{banxiaosi}{13bp}
\nju@def@fontsize{xiaosi}{12bp}
\nju@def@fontsize{dawu}{11bp}
\nju@def@fontsize{wuhao}{10.5bp}
\nju@def@fontsize{xiaowu}{9bp}
\nju@def@fontsize{liuhao}{7.5bp}
\nju@def@fontsize{xiaoliu}{6.5bp}
\nju@def@fontsize{qihao}{5.5bp}
\nju@def@fontsize{bahao}{5bp}
%    \end{macrocode}
% \end{macro}
% \end{macro}
% \end{macro}
% \end{macro}
% \end{macro}
% \end{macro}
% \end{macro}
% \end{macro}
% \end{macro}
% \end{macro}
% \end{macro}
% \end{macro}
% \end{macro}
% \end{macro}
% \end{macro}
% \end{macro}
% \end{macro}
% \end{macro}
%
%
% \subsection{中文标点}
%
% \newcommand\unicodechar[1]{U+#1（\symbol{"#1}）}
% 由于 Unicode 的一些标点符号是中西文混用的：
% \unicodechar{00B7}、
% \unicodechar{2013}、
% \unicodechar{2014}、
% \unicodechar{2018}、
% \unicodechar{2019}、
% \unicodechar{201C}、
% \unicodechar{201D}、
% \unicodechar{2025}、
% \unicodechar{2026}、
% \unicodechar{2E3A}，
% 所以要根据语言设置正确的字体。
% \footnote{\url{https://github.com/CTeX-org/ctex-kit/issues/389}}
% 所以要根据语言设置正确的字体。
%    \begin{macrocode}
\newcommand\nju@setchinese{%
  \xeCJKResetPunctClass
}
\newcommand\nju@setenglish{%
  \xeCJKDeclareCharClass{HalfLeft}{"2018, "201C}%
  \xeCJKDeclareCharClass{HalfRight}{
    "00B7, "2019, "201D, "2013, "2014, "2025, "2026, "2E3A,
  }%
}
\newcommand\nju@setdefaultlanguage{%
  \ifnju@chinese
    \nju@setchinese
  \else
    \nju@setenglish
  \fi
}
%    \end{macrocode}
%
% \subsection{局部设置}
% \subsubsection{页眉页脚}
% \label{sec:headerfooter}
%
% 定义页眉和页脚样式。
% \begin{macro}{\ps@nju@empty}
% \begin{macro}{\ps@nju@plain}
% \begin{macro}{\ps@nju@headings}
% \begin{itemize}
% \item \texttt{nju@empty}：页眉页脚都没有
% \item \texttt{nju@plain}：只显示页脚的页码。\cs{chapter} 自动调用
% \cs{thispagestyle\{nju@plain\}}。
% \item \texttt{nju@headings}：页眉页脚同时显示
% \end{itemize}
%    \begin{macrocode}
\fancypagestyle{nju@empty}{%
  \fancyhf{}
  \renewcommand{\headrulewidth}{0pt}
  \renewcommand{\footrulewidth}{0pt}}
\fancypagestyle{nju@plain}{%
  \fancyhead{}
  \fancyfoot[C]{\xiaowu\thepage}
  \renewcommand{\headrulewidth}{0pt}
  \renewcommand{\footrulewidth}{0pt}}
\fancypagestyle{nju@headings}{%
  \fancyhead{}
  \fancyhead[C]{\wuhao\normalfont\leftmark}
  \fancyfoot{}
  \fancyfoot[C]{\wuhao\thepage}
  \renewcommand{\headrulewidth}{0.4pt}
  \renewcommand{\footrulewidth}{0pt}}
%    \end{macrocode}
% \end{macro}
% \end{macro}
% \end{macro}
%
% \subsubsection{段落}
% \label{sec:paragraph}
%
% 全文首行缩进 2 字符，标点符号用全角
%    \begin{macrocode}
\ctexset{%
  punct=quanjiao,
  space=auto,
  autoindent=true}
%    \end{macrocode}
%
% \subsubsection{列表}
% 利用 \pkg{enumitem} 命令调整默认列表环境间的距离，以符合中文习惯。
%    \begin{macrocode}
\RequirePackage[shortlabels]{enumitem}
\RequirePackage{environ}
\setlist{nosep}
%    \end{macrocode}
%
%
% \subsubsection{脚注}
% \label{sec:footnote}
% 脚注符合中文习惯，数字带圈。
%    \begin{macrocode}
\ifthenelse{\equal{\nju@fontset}{mac}}{
  \newfontfamily\nju@circlefont{Songti SC Light}
}{
  \ifthenelse{\equal{\nju@fontset}{windows}}{
    \newfontfamily\nju@circlefont{SimSun}
  }{
    \IfFontExistsTF{XITS-Regular.otf}{
      \newfontfamily\nju@circlefont{XITS-Regular.otf}
    }{
      \newfontfamily\nju@circlefont{xits-regular.otf}
    }
  }
}
\def\nju@textcircled#1{%
  \ifnum\value{#1} >9%
    \ClassError{njurepo}%
      {Too many footnotes in this page.}{Keep footnote less than 10.}%
  \fi
  {\nju@circlefont\symbol{\numexpr\value{#1}+"245F\relax}}%
}
\renewcommand{\thefootnote}{\nju@textcircled{footnote}}
\renewcommand{\thempfootnote}{\nju@textcircled{mpfootnote}}
%    \end{macrocode}
%
% 定义脚注分割线，字号（宋体小五），以及悬挂缩进（1.5字符）。
%    \begin{macrocode}
\def\footnoterule{\vskip-3\p@\hrule\@width0.3\textwidth\@height0.4\p@\vskip2.6\p@}
\let\nju@footnotesize\footnotesize
\renewcommand\footnotesize{\nju@footnotesize\xiaowu[1.5]}
%\footnotemargin1.5em\relax
%    \end{macrocode}
%
% \cs{@makefnmark} 默认是上标样式，而在脚注部分要求为正文大小。利用\cs{patchcmd} 动态调整 \cs{@makefnmark} 的定义。
%    \begin{macrocode}
\let\nju@makefnmark\@makefnmark
\def\nju@@makefnmark{\hbox{{\normalfont\@thefnmark}}}
\pretocmd{\@makefntext}{\let\@makefnmark\nju@@makefnmark}{}{}
\apptocmd{\@makefntext}{\let\@makefnmark\nju@makefnmark}{}{}
%    \end{macrocode}
%
%
% \subsubsection{定理环境}
% \label{sec:equation}
%
% 定理标题使用黑体，正文使用宋体，冒号隔开。
%    \begin{macrocode}
\theorembodyfont{\normalfont}
\theoremheaderfont{\normalfont\heiti}
\theoremsymbol{\ensuremath{\square}}
\newtheorem*{proof}{证明}
\theoremstyle{plain}
\theoremsymbol{}
\theoremseparator{：}
\ifnju@chinese
  \newcommand\nju@assumption@name{假设}
  \newcommand\nju@definition@name{定义}
  \newcommand\nju@proposition@name{命题}
  \newcommand\nju@lemma@name{引理}
  \newcommand\nju@theorem@name{定理}
  \newcommand\nju@axiom@name{公理}
  \newcommand\nju@corollary@name{推论}
  \newcommand\nju@exercise@name{练习}
  \newcommand\nju@example@name{例}
  \newcommand\nju@remark@name{注释}
  \newcommand\nju@problem@name{问题}
  \newcommand\nju@conjecture@name{猜想}
  \newcommand\nju@solution@name{解}
\else
  \newcommand\nju@assumption@name{Assumption}
  \newcommand\nju@definition@name{Definition}
  \newcommand\nju@proposition@name{Proposition}
  \newcommand\nju@lemma@name{Lemma}
  \newcommand\nju@theorem@name{Theorem}
  \newcommand\nju@axiom@name{Axiom}
  \newcommand\nju@corollary@name{Corollary}
  \newcommand\nju@exercise@name{Exercise}
  \newcommand\nju@example@name{Example}
  \newcommand\nju@remark@name{Remark}
  \newcommand\nju@problem@name{Problem}
  \newcommand\nju@conjecture@name{Conjecture}
  \newcommand\nju@solution@name{Solution}
\fi
\theoremheaderfont{\bfseries}
\newtheorem{assumption}{\nju@assumption@name}[chapter]
\newtheorem{definition}{\nju@definition@name}[chapter]
\newtheorem{proposition}{\nju@proposition@name}[chapter]
\newtheorem{lemma}{\nju@lemma@name}[chapter]
\newtheorem{theorem}{\nju@theorem@name}[chapter]
\newtheorem{axiom}{\nju@axiom@name}[chapter]
\newtheorem{corollary}{\nju@corollary@name}[chapter]
\newtheorem{exercise}{\nju@exercise@name}[chapter]
\newtheorem{example}{\nju@example@name}[chapter]
\newtheorem{remark}{\nju@remark@name}[chapter]
\newtheorem{problem}{\nju@problem@name}[chapter]
\newtheorem{conjecture}{\nju@conjecture@name}[chapter]
\newtheorem{solution}{\nju@solution@name}[chapter]

%\RequirePackage{microtype}
\ifnju@chinese
\newcommand{\promisewords}{请独立完成作业，不得抄袭。\\若参考了其它资料，请给出引用。\\鼓励讨论，但需独立书写解题过程。}
\else
\newcommand{\promisewords}{I promise this work is done on my own with no plagiarism.}
\fi
\newcommand{\pshw}{\section*{\scshape Part I\ \ \ Homework}}
\newcommand{\pscr}{\section*{\scshape Part II\ \ \ Correction}}
\newcommand{\psfb}{\section*{\scshape Part III\ \ \ Feedback}}
\newcommand{\Hrule}{\noindent\rule{\linewidth}{0.5mm}}

\ifnju@awesomefont
\RequirePackage{awesomefont}
\fi

\theorempostwork{\vspace{-0.5cm}\Hrule}
\newtheorem*{pssolution}{\ifnju@awesomefont\faPencilSquareO\ \fi\nju@solution@name}
\RequirePackage[listings]{tcolorbox}
\newtcolorbox{ps@problem}[1]{fonttitle=\bfseries,title=#1,before skip=0.5cm, after skip=-0.5cm}
\newenvironment{psproblem}[1][]{
    \begin{ps@problem}{\ifnju@awesomefont\faQuestionCircle\ \fi\nju@problem@name\ #1}
}{
    \end{ps@problem}
}
%
% \subsubsection{浮动对象}
% \label{sec:float}
% 设置浮动对象和文字之间的距离
%    \begin{macrocode}
\setlength{\floatsep}{20bp \@plus4pt \@minus1pt}
\setlength{\intextsep}{20bp \@plus4pt \@minus2pt}
\setlength{\textfloatsep}{20bp \@plus4pt \@minus2pt}
\setlength{\@fptop}{0bp \@plus1.0fil}
\setlength{\@fpsep}{12bp \@plus2.0fil}
\setlength{\@fpbot}{0bp \@plus1.0fil}
%    \end{macrocode}
%
% 下面这组命令使浮动对象的缺省值稍微宽松一点，从而防止幅度对象占据过多的文本页面，
% 也可以防止在很大空白的浮动页上放置很小的图形。
%    \begin{macrocode}
\renewcommand{\textfraction}{0.15}
\renewcommand{\topfraction}{0.85}
\renewcommand{\bottomfraction}{0.65}
\renewcommand{\floatpagefraction}{0.60}
%    \end{macrocode}
%
% 定制浮动图形和表格标题样式
% \begin{itemize}
%   \item 图表标题字体为 11pt， 这里写作大五号
%   \item 去掉图表号后面的冒号。图序与图名文字之间空一个汉字符宽度。
%   \item 图：caption 在下，段前空 6 磅，段后空 12 磅
%   \item 表：caption 在上，段前空 12 磅，段后空 6 磅
% \end{itemize}
%
%    \begin{macrocode}
\let\old@tabular\@tabular
\def\nju@tabular{\dawu[1.5]\old@tabular}
\DeclareCaptionLabelFormat{nju}{{\dawu[1.5]\normalfont #1~#2}}
\DeclareCaptionLabelSeparator{nju}{\hspace{1em}}
\DeclareCaptionFont{nju}{\dawu[1.5]}
\captionsetup{labelformat=nju,labelsep=nju,font=nju,skip=6bp}
\captionsetup[table]{position=top}
\captionsetup[figure]{position=bottom}
\captionsetup[sub]{font=nju}
\renewcommand{\thesubfigure}{(\alph{subfigure})}
\renewcommand{\thesubtable}{(\alph{subtable})}
% \renewcommand{\p@subfigure}{:}
%    \end{macrocode}
% 我们采用 \pkg{longtable} 来处理跨页的表格。同样我们需要设置其默认字体为五号。
%    \begin{macrocode}
\let\nju@LT@array\LT@array
\def\LT@array{\dawu[1.5]\nju@LT@array} % set default font size
%    \end{macrocode}
%
% \begin{macro}{\hlinewd}
% 简单的表格使用三线表推荐用 \cs{hlinewd}。如果表格比较复杂还是用 \pkg{booktabs} 的命令好一些。
%    \begin{macrocode}
\def\hlinewd#1{%
  \noalign{\ifnum0=`}\fi\hrule \@height #1 \futurelet
    \reserved@a\@xhline}
%    \end{macrocode}
% \end{macro}
%
%
% \subsubsection{章节标题}
% \label{sec:theor}
%    \begin{macrocode}
\ifnju@chinese
  \ctexset{%
    chapter/name={第,章},
    appendixname=附录,
    contentsname={目\hspace{\ccwd}录},
    listfigurename=插图索引,
    listtablename=表格索引,
    figurename=图,
    tablename=表,
    bibname=参考文献,
    indexname=索引,
  }
  \newcommand\listequationname{公式索引}
  \newcommand\equationname{公式}
\else
  \newcommand\listequationname{List of Equations}
  \newcommand\equationname{Equation}
\fi
\newcommand{\cabstractname}{摘\hspace{\ccwd}要}
\newcommand{\eabstractname}{Abstract}
\let\CJK@todaysave=\today
\def\CJK@todaysmall@short{\the\year 年 \the\month 月}
\def\CJK@todaysmall{\the\year 年 \the\month 月 \the\day 日}
\def\CJK@todaybig@short{\zhdigits{\the\year}年\zhnumber{\the\month}月}
\def\CJK@todaybig{\zhdigits{\the\year}年\zhnumber{\the\month}月\zhnumber{\the\day}日}
\def\CJK@today{\CJK@todaysmall}
\renewcommand\today{\CJK@today}
\newcommand\CJKtoday[1][1]{%
  \ifcase#1\def\CJK@today{\CJK@todaysave}
    \or\def\CJK@today{\CJK@todaysmall}
    \or\def\CJK@today{\CJK@todaybig}
  \fi}
%    \end{macrocode}
%
% \pkg{fancyhdr} 定义页眉页脚很方便，但是有一个非常隐蔽的坑。通过 \pkg{fancyhdr}
% 定义的样式在第一次被调用时会修改 \cs{chaptermark}，这会导致页眉信息错误（多余
% 章号并且英文大写）。这是因为在原始的 \file{book.cls} 中定义如下（大意）：
% \begin{latex}
% \newcommand\chaptername{Chapter}
% \newcommand\@chapapp{\chaptername}
% \def\chaptermark#1{
%   \markboth{\MakeUppercase{\@chapapp\ \thechapter}}{}}
% \end{latex}
% 很显然这个 \cs{\@chapapp} 不适合中文，因此我们使用\cs{CTEXthechapter}（
% 如，“第 x 章”），同时会将 \cs{MakeUppercase} 去掉。也就是说我们会做如下动作：
% \begin{latex}
% \renewcommand{\chaptermark}[1]{\@mkboth{\CTEXthechapter\hskip\ccwd#1}{}}
% \end{latex}
% 但，\pkg{fancyhdr} 不知何故在 \cs{ps@fancy} 中对 \cs{chaptermark} 进行重定义
% （其实一模一样），而这个 \cs{ps@fancy} 会在 \cs{fancypagestyle} 中使用，如下：
% \begin{latex}
% \newcommand{\fancypagestyle}[2]{%
%   \@namedef{ps@#1}{\let\fancy@gbl\relax#2\relax\ps@fancy}}
% \end{latex}
% 这样的话，\cs{ps@fancy} 会在 \pkg{fancyhdr} 定义的任何样式首次样被激活时调用，从
% 而覆盖我们的 \cs{chaptermark} 定义（后续样式再激活不会重复覆盖）。所以我们采用如下
% 方法解决：
%    \begin{macrocode}
\AtBeginDocument{%
  \pagestyle{nju@empty}
  \renewcommand{\chaptermark}[1]{\@mkboth{\CTEXthechapter\hskip\ccwd#1}{}}}
%    \end{macrocode}
%
% 各级标题格式设置。
% \begin{description}
% \item[chapter] 章序号与章名之间空一个汉字符 黑体三号字，居中书写，单倍行距，段
%   前空 24 磅，段后空 18 磅。本科要求：段前段后间距 30/20 pt，行距 20pt。但正文
%   章节 30pt 的话和样例效果不一致。
% \item[section] 一级节标题，例如：\fbox{2.1 实验装置与实验方法}。节标题序号与标
%   题名之间空一个汉字符（下同）。采用黑体四号（14pt）字居左书写，行距为固定
%   值 20 磅，段前空 24 磅，段后空 6 磅。本科：25/12 pt，行距 18pt。
% \item[subsection] 二级节标题，例如：\fbox{2.1.1 实验装置}。采用黑体 13pt 字居左
%   书写，行距为固定值 20 磅，段前空 12 磅，段后空 6 磅。本科：中文黑体 12pt 字，
%   英文 13pt 字，段间距 12/6 pt，行距 15pt。
% \item[subsubsection] 三级节标题，例如：\fbox{2.1.2.1 归纳法}。采用黑体小四号
%   （12pt）字居左书写，行距为固定值 20 磅，段前空 12 磅，段后空 6 磅。
%
% \end{description}
%    \begin{macrocode}
\newcommand\nju@chapter@titleformat[1]{%
    \ifthenelse%
      {\equal{#1}{\eabstractname}}%
      {\bfseries #1}%
      {#1}%
  }
\ctexset{%
  chapter={
    afterindent=true,
    pagestyle={nju@headings},
    beforeskip={9bp},
    aftername=\hskip\ccwd,
    afterskip={24bp},
    format={\centering\sffamily\sanhao[1]},
    nameformat=\relax,
    numberformat=\relax,
    titleformat=\nju@chapter@titleformat,
    lofskip=0pt,
    lotskip=0pt,
  },
  section={
    afterindent=true,
    beforeskip={24bp\@plus 1ex \@minus .2ex},
    afterskip={6bp\@plus .2ex},
    format={\sffamily\sihao[1.429]},
  },
  subsection={
    afterindent=true,
    beforeskip={16bp\@plus 1ex \@minus .2ex},
    afterskip={6bp \@plus .2ex},
    format={\sffamily\banxiaosi[1.538]},
    numberformat={\sffamily\banxiaosi[1.538]},
  },
  subsubsection={
    afterindent=true,
    beforeskip={16bp\@plus 1ex \@minus .2ex},
    afterskip={6bp \@plus .2ex},
    format={\sffamily\xiaosi[1.667]},
  },
  paragraph/afterindent=true,
  subparagraph/afterindent=true}
%    \end{macrocode}
%
% \begin{macro}{\nju@chapter*}
% 默认的 \cs{chapter*} 很难同时满足研究生院和本科生的论文要求。本科论文要求所有的
% 章都出现在目录里，比如摘要、Abstract、主要符号表等，所以可以简单的扩展默
% 认\cs{chapter*} 实现这个目的。但是研究生又不要这些出现在目录中，而且致谢和声明
% 部分的章名、页眉和目录都不同，所以定义一个灵活的 \cs{nju@chapter*} 专门处理这些
% 要求。
%
% \cs{nju@chapter*}\oarg{tocline}\marg{title}\oarg{header}: tocline 是出现在目录
% 中的条目，如果为空则此 chapter 不出现在目录中，如果省略表示目录出现 title；
% title 是章标题；header 是页眉出现的标题，如果忽略则取 title。通过这个宏我才真
% 正体会到 \TeX\ macro 的力量！
%    \begin{macrocode}
\newcounter{nju@bookmark}
\NewDocumentCommand\nju@chapter{s o m o}{
  \IfBooleanF{#1}{%
    \ClassError{njurepo}{You have to use the star form: \string\nju@chapter*}{}
  }%
  \ifnju@right\cleardoublepage\else\clearpage\fi\phantomsection%
  \IfValueTF{#2}{%
    \ifthenelse{\equal{#2}{}}{%
      \addtocounter{nju@bookmark}\@ne
      \pdfbookmark[0]{#3}{njuchapter.\thenju@bookmark}
    }{%
      \addcontentsline{toc}{chapter}{#3}
    }
  }{%
    \addcontentsline{toc}{chapter}{#3}
  }%
  \ctexset{chapter/beforeskip=25bp}
  \chapter*{#3}%
  \ctexset{chapter/beforeskip=15bp}
  \IfValueTF{#4}{%
    \ifthenelse{\equal{#4}{}}
    {\@mkboth{}{}}
    {\@mkboth{#4}{#4}}
  }{%
    \@mkboth{#3}{#3}
  }
}
%    \end{macrocode}
% \end{macro}
%
%
% \subsubsection{目录}
% \label{sec:toc}
% 最多 4 层，即: x.x.x.x，对应的命令和层序号分别是：
% \cs{chapter}(0), \cs{section}(1), \cs{subsection}(2), \cs{subsubsection}(3)。
%    \begin{macrocode}
\setcounter{secnumdepth}{3}
\setcounter{tocdepth}{2}
%    \end{macrocode}
%
% 每章标题行前空 6 磅，后空 0 磅。章节名中英文用 Arial 字体，页码仍用 Times。
% \begin{macro}{\tableofcontents}
%    \begin{macrocode}
\renewcommand\tableofcontents{%
  \nju@chapter*[]{\contentsname}
  \xiaosi[1.65]\@starttoc{toc}\normalsize}
%    \end{macrocode}
% 调整目录样式
%    \begin{macrocode}
\def\@pnumwidth{2em}
\def\@tocrmarg{\@pnumwidth}
\def\@dotsep{1}
\renewcommand*\l@chapter[2]{%
  \ifnum \c@tocdepth >\m@ne
    \addpenalty{-\@highpenalty}%
    \vskip 4bp \@plus\p@
    \setlength\@tempdima{4em}%
    \begingroup
      \parindent \z@ \rightskip \@pnumwidth
      \parfillskip -\@pnumwidth
      \leavevmode
      \advance\leftskip\@tempdima
      \hskip -\leftskip
      {#1}%
      \leaders\hbox{$\m@th\mkern \@dotsep mu\hbox{.}\mkern \@dotsep mu$}\hfill%
      \nobreak{#2}\par
      \penalty\@highpenalty
    \endgroup
  \fi}

\patchcmd{\@dottedtocline}{\hb@xt@\@pnumwidth}{\hbox}{}{}
\renewcommand*\l@section{%
  \@dottedtocline{1}{\ccwd}{2.1em}}
\renewcommand*\l@subsection{%
  \@dottedtocline{2}{2\ccwd}{3em}}
\renewcommand*\l@subsubsection{%
  \@dottedtocline{3}{3.5em}{3.8em}}
%    \end{macrocode}
% \end{macro}
%
% \subsection{附加页面}
% \label{sec:etc}
%
% \subsubsection{封面}
% \label{sec:cover}
% 定义封面参数。
%    \begin{macrocode}
\def\nju@def@term#1{%
  \define@key{nju}{#1}{\csname #1\endcsname{##1}}
  \expandafter\gdef\csname #1\endcsname##1{%
    \expandafter\gdef\csname nju@#1\endcsname{##1}}
  \csname #1\endcsname{}}
\nju@def@term{ctitle}
\nju@def@term{csubtitle}
\nju@def@term{csubsubtitle}
\nju@def@term{etitle}
\nju@def@term{esubtitle}
\nju@def@term{esubsubtitle}
\nju@def@term{cauthor}
\nju@def@term{csupervisor}
\nju@def@term{cassosupervisor}
\nju@def@term{ccosupervisor}
\nju@def@term{eauthor}
\nju@def@term{esupervisor}
\nju@def@term{eassosupervisor}
\nju@def@term{ecosupervisor}
\nju@def@term{cdegree}
\nju@def@term{edegree}
\nju@def@term{cdepartment}
\nju@def@term{edepartment}
\nju@def@term{cmajor}
\nju@def@term{emajor}
\nju@def@term{cdate}
\nju@def@term{edate}
\nju@def@term{stdid}
\nju@def@term{mail}
\cdate{\CJK@todaybig@short}
\edate{\ifcase \month \or January\or February\or March\or April\or May%
       \or June\or July \or August\or September\or October\or November
       \or December\fi\unskip,\ \ \the\year}
%    \end{macrocode}
%
% \begin{environment}{cabstract}
% \begin{environment}{eabstract}
% 摘要最好以环境的形式出现（否则命令的形式会导致开始结束的括号距离太远，我不喜
% 欢），这就必须让环境能够自己保存内容留待以后使用。使用 \pkg{environ} 的
% \cs{Collect@Body} 来实现。
%    \begin{macrocode}
\newcommand{\nju@@cabstract}[1]{\long\gdef\nju@cabstract{#1}}
\newenvironment{cabstract}{\Collect@Body\nju@@cabstract}{}
\newcommand{\nju@@eabstract}[1]{\long\gdef\nju@eabstract{#1}}
\newenvironment{eabstract}{\Collect@Body\nju@@eabstract}{}
%    \end{macrocode}
% \end{environment}
% \end{environment}
%
% \begin{macro}{\nju@parse@keywords}
%   不同论文格式关键词之间的分割不太相同，我们用 \cs{ckeywords} 和
%    \cs{ekeywords} 来收集关键词列表，然后用本命令来生成符合要求的格式。
%    \begin{macrocode}
\def\nju@parse@keywords#1{
  \define@key{nju}{#1}{\csname #1\endcsname{##1}}
  \expandafter\gdef\csname nju@#1\endcsname{}
  \expandafter\gdef\csname #1\endcsname##1{
    \@for\reserved@a:=##1\do{
      \expandafter\ifx\csname nju@#1\endcsname\@empty\else
        \expandafter\g@addto@macro\csname nju@#1\endcsname{%
          \ignorespaces\csname nju@#1@separator\endcsname}
      \fi
      \expandafter\expandafter\expandafter\g@addto@macro%
        \expandafter\csname nju@#1\expandafter\endcsname\expandafter{\reserved@a}}}}
%    \end{macrocode}
% \end{macro}
% \begin{macro}{\ckeywords}
% \begin{macro}{\ekeywords}
% 利用 \cs{nju@parse@keywords} 来定义，内部通过 \cs{nju@ckeywords} 和
% \cs{nju@ekeywords} 来引用。
% \changes{v3.1}{2007/06/16}{增强的关键词命令。}
%    \begin{macrocode}
\nju@parse@keywords{ckeywords}
\nju@parse@keywords{ekeywords}
%    \end{macrocode}
% \end{macro}
% \end{macro}
%
% \begin{macro}{\njusetup}
% 由上可见，封面和封底有一大堆信息需要设置，为了简化操作界面，提供一
% 个 \cs{njusetup} 命令支持 key/value 的方式来设置。key 就是前面各个设置项的
% 名字。\note[说明：]{只能设置普通项，不支持环境项，
% 如 \texttt{cabstract} 和 \texttt{eabstract}。} 由于这些设置项被 \cs{makecover}
% 调用，所以此命令需要在 \cs{makecover} 之前被调用。
%    \begin{macrocode}
\def\njusetup{\kvsetkeys{nju}}
%    \end{macrocode}
% \end{macro}
%
% 定义封面用到的各种文字。
%    \begin{macrocode}
\def\nju@ckeywords@separator{；}
\def\nju@ekeywords@separator{;}
\def\nju@catalog@number@title{分类号}
\def\nju@id@title{编号}
\def\nju@title@sep{：}
\def\nju@schoolname{南京大学}
\def\nju@author@title{姓名}
\def\nju@department@title{系别}
\def\nju@major@title{专业}
\def\nju@supervisor@title{指导教师}
\def\nju@assosuper@title{辅导教师}
\def\nju@studentid@title{学号}
\def\nju@date@title{日期}
\def\nju@mail@title{邮箱}
\newcommand{\nju@ckeywords@title}{关键词：}
\def\nju@title@pre{}

\def\nju@eng@title@sep{:}
\def\nju@eng@author@title{Name}
\def\nju@eng@studentid@title{StdID}
\def\nju@eng@date@title{Date}
\def\nju@eng@mail@title{E-mail}
%    \end{macrocode}
%
% 中文小型标题
%    \begin{macrocode}
\renewcommand{\maketitle}{
  \nju@setup@pdfinfo
  \begin{center} {\LARGE \ifnju@chinese\nju@ctitle\else\nju@etitle\fi}
  \end{center}
  \hspace*{\fill}
  \ifnju@chinese
    \nju@author@title\nju@title@sep\CJKunderline{\nju@cauthor}
  \else
    \nju@eng@author@title\nju@eng@title@sep\underline{\nju@eauthor}
  \fi
  \hspace*{\fill}
  \ifx\nju@stdid\@empty\relax
  \else
    \ifnju@chinese
      \nju@studentid@title\nju@title@sep\CJKunderline{\nju@stdid}
    \else
      \nju@eng@studentid@title\nju@eng@title@sep\underline{\nju@stdid}
    \fi
  \fi
  \hspace*{\fill}
  \ifnju@chinese
    \nju@date@title\nju@title@sep\CJKunderline{\today}
  \else
    \nju@eng@date@title\nju@eng@title@sep\CJKunderline{\nju@edate}
  \fi
  \hspace*{\fill}\\
}
%    \end{macrocode}
%
% 别样封面
%    \begin{macrocode}
\newcommand{\maketitlepage}{
  \nju@setup@pdfinfo
  \begin{titlepage}
    \begin{center}
    \ifx\nju@esubsubtitle\@empty\relax  {\LARGE\sffamily\scshape\ifnju@chinese\nju@csubsubtitle\else\nju@esubsubtitle\fi\ }\\[1.5cm]
    \else
    {\LARGE\sffamily\scshape \ifnju@chinese\nju@csubsubtitle\else\nju@esubsubtitle\fi}\\[1.5cm]
    \fi
		{\Large\sffamily\scshape \ifnju@chinese\nju@csubtitle\else\nju@esubtitle\fi}\\
    \rule{\linewidth}{0.5mm} \\[0.4cm]
		{\huge\sffamily\bfseries \ifnju@chinese\nju@ctitle\else\nju@etitle\fi}\\
		\rule{\linewidth}{0.5mm} \\[1.5cm]
			
		\begin{center}
			\begin{tabular}{@{\hspace{0.5cm}}l@{\hspace{0.5cm}}l}
				\nju@eauthor & \nju@stdid\\
			\end{tabular}
		\end{center}
		\vfill
		{\large \nju@edate}
    \end{center}
    \ifnju@right\cleardoublepage\else\clearpage\fi
  \end{titlepage}
}
%    \end{macrocode}
%
% \myentry{封面第一页}
% \begin{macro}{\nju@first@titlepage}
% 题名使用一号黑体字，一行写不下时可分两行写，并采用 1.25 倍行距。
% 申请学位的学科门类: 小二号宋体字。
% 中文封面页边距：
%  上- 6.0 厘米，下- 5.5 厘米，左- 4.0 厘米，右- 4.0 厘米，装订线 0 厘米；
%
%    \begin{macrocode}
\newcommand\nju@underline[2][6em]{\hskip1pt\underline{\hb@xt@ #1{\hss#2\hss}}\hskip3pt}
\newlength{\nju@title@width}
\ifxetex % todo: ugly codes
  \newcommand{\nju@put@title}[2][\nju@title@width]{%
  \begin{CJKfilltwosides}[b]{#1}#2\end{CJKfilltwosides}}
\else
  \newcommand{\nju@put@title}[2][\nju@title@width]{%
  \begin{CJKfilltwosides}{#1}#2\end{CJKfilltwosides}}
\fi
\newcommand{\nju@first@titlepage}{
  \begin{center}
    \vspace*{-1.6cm}
    \parbox[b][2.4cm][t]{\textwidth}{%
      \rule{1cm}{0cm}}
      \vskip0.65cm
      {\includegraphics[width=0.3\textwidth]{njuname0.pdf}}
      \par\vskip1.5cm
      {\xiaochu\heiti\ziju{0.5}\textbf\nju@csubtitle}
      \vskip2.2cm\hskip0.8cm
      \noindent\heiti\xiaoer\nju@title@pre
      \parbox[t]{12cm}{%
      \ignorespaces\yihao[1.51]%
      \renewcommand{\CJKunderlinebasesep}{0.25cm}%
      \renewcommand{\ULthickness}{1.3pt}%
      \ifxetex
        \xeCJKsetup{underline/format=\color{black}}%
      \else
        \def\CJKunderlinecolor{\color{black}}%
      \fi
      \centering\CJKunderline*{\nju@ctitle}
      
    }%
      \vskip1.3cm
%    \end{macrocode}
%
% 作者及导师信息部分使用三号仿宋字
%    \begin{macrocode}
      \vskip0.75cm
      \ifx\nju@cassosupervisor\@empty%
        \def\nju@tempa{7.15cm}
      \else%
        \def\nju@tempa{8.15cm}
      \fi%
      \parbox[t][\nju@tempa][t]{\textwidth}{%
        {\fangsong\sanhao[1.95]%
         \hspace*{1.9cm}
         \setlength{\nju@title@width}{4em}
         \setlength{\extrarowheight}{6pt}
         \ifxetex % todo: ugly codes
           \begin{tabular}{p{\nju@title@width}@{}l@{\extracolsep{8pt}}l}
         \else
           \begin{tabular}{p{\nju@title@width}l@{}l}
         \fi
             \nju@put@title{\nju@department@title} & \nju@title@sep
               & \nju@cdepartment\\
             \nju@put@title{\nju@major@title}      & \nju@title@sep
               & \nju@cmajor\\
             \nju@put@title{\nju@author@title}     & \nju@title@sep
               & \nju@cauthor \\
             \nju@put@title{\nju@supervisor@title} & \nju@title@sep
               & \nju@csupervisor\\
             \ifx\nju@cassosupervisor\@empty\else%
               \nju@put@title{\nju@assosuper@title} & \nju@title@sep
               & \nju@cassosupervisor\\
             \fi
           \end{tabular}
        }}
%    \end{macrocode}
%
% 论文成文打印的日期，用三号宋体汉字，不用阿拉伯数字
% 本科：论文成文打印的日期用阿拉伯数字，采用小四号宋体
%    \begin{macrocode}
     \begin{center}
       {\vskip-1.0cm\xiaosi
         \songti\nju@cdate}
     \end{center}
    \end{center}} % end of titlepage
%    \end{macrocode}
% \end{macro}
%
% \myentry{英文封面}
% \begin{macro}{\nju@engcover}
%    \begin{macrocode}
\newcommand{\nju@engcover}{%
  \begin{center}
    \vspace*{-5pt}
    \parbox[t][5.2cm][t]{\paperwidth-7.2cm}{
      \renewcommand{\baselinestretch}{1.5}
      \begin{center}
        \erhao[1.1]\bfseries\sffamily\nju@etitle%
      \end{center}}
    \parbox[t][][b]{\paperwidth-7.2cm}{
      \renewcommand{\baselinestretch}{1.3}
      \begin{center}
        \sanhao\sffamily by\\[3bp]
        \bfseries\nju@eauthor%
        \ifx\nju@emajor\empty\relax\else
          \\(~\nju@emajor~)%
        \fi
      \end{center}}
    \par\vspace{0.9cm}
    \parbox[t][2.1cm][t]{\paperwidth-7.2cm}{
      \renewcommand{\baselinestretch}{1.2}
      \xiaosan\centering
      \begin{tabular}{rl}
        Supervisor : & \nju@esupervisor\\
        \ifx\nju@eassosupervisor\@empty
          \else Associate Supervisor : & \nju@eassosupervisor\\\fi
        \ifx\nju@ecosupervisor\@empty
          \else Cooperate Supervisor : & \nju@ecosupervisor\\\fi
      \end{tabular}}
    \parbox[t][2cm][b]{\paperwidth-7.2cm}{
    \begin{center}
      \sanhao\bfseries\sffamily\nju@edate
    \end{center}}
  \end{center}}
%    \end{macrocode}
% \end{macro}
%
% \begin{macro}{\makecover}
% 生成封面总命令。
%    \begin{macrocode}
\def\makecover{%
  \nju@setup@pdfinfo\nju@makecover}
\def\nju@setup@pdfinfo{%
  \ifnju@chinese
    \hypersetup{
      pdftitle    = \nju@ctitle,
      pdfauthor   = \nju@cauthor,
      pdfsubject  = \nju@cdegree,
      pdfkeywords = \nju@ckeywords,
    }%
  \else
    \hypersetup{
      pdftitle    = \nju@etitle,
      pdfauthor   = \nju@eauthor,
      pdfsubject  = \nju@edegree,
      pdfkeywords = \nju@ekeywords,
    }%
  \fi
  \hypersetup{
    pdfcreator={\njurepo-v\version}}}
\NewDocumentCommand{\nju@makecover}{o}{
  \phantomsection
  \pdfbookmark[-1]{\nju@ctitle}{ctitle}
  \normalsize%
  \begin{titlepage}
    \ifnju@chinese
      \nju@first@titlepage
    \else
      \nju@engcover
    \fi
    \ifnju@right\cleardoublepage\else\clearpage\fi
  \end{titlepage}
}
\newcommand{\makeabstract}{
  \normalsize
  \nju@makeabstract
  \let\@tabular\nju@tabular
}
%    \end{macrocode}
% \end{macro}
%
% \subsubsection{摘要}
% \label{sec:abstractformat}
%
% \begin{macro}{\nju@put@keywords}
% 排版关键字。
%    \begin{macrocode}
\newbox\nju@kw
\newcommand\nju@put@keywords[2]{%
  \begingroup
    \setbox\nju@kw=\hbox{#1}
    \indent%
    \box\nju@kw#2\par
  \endgroup}
%    \end{macrocode}
% \end{macro}
%
% \begin{macro}{\nju@makeabstract}
% 中文摘要部分的标题为“\textbf{摘要}”，用黑体三号字。
%    \begin{macrocode}
\newcommand{\nju@makeabstract}{%
  \clearpage
  \pagestyle{nju@plain}
  \pagenumbering{Roman}
%    \end{macrocode}
%
% 摘要内容用小四号字书写，两端对齐，汉字用宋体，外文字用 Times New Roman 体，
% 标点符号一律用中文输入状态下的标点符号。
%    \begin{macrocode}
  \ifnju@chinese
    \nju@setchinese
    \nju@chapter*[]{\cabstractname} % no tocline
    \nju@cabstract
    \vskip12bp
    \nju@put@keywords{\textbf\nju@ckeywords@title}{\nju@ckeywords}
  \else
  \nju@setenglish
    \nju@chapter*[]{\eabstractname} % no tocline
    \nju@eabstract
    \vskip12bp
    \nju@put@keywords{%
      \textbf{Key Words:\enskip}}{\nju@ekeywords}%
  \fi
  \nju@setdefaultlanguage
}
%    \end{macrocode}
% \end{macro}
%
%
%
% \subsubsection{主要符号表}
% \label{sec:denotationfmt}
% \begin{environment}{denotation}
% 主要符号对照表。
%    \begin{macrocode}
\ifnju@chinese
  \newcommand\nju@denotation@name{主要符号对照表}
\else
  \newcommand\nju@denotation@name{Nomenclature}
\fi
\newenvironment{denotation}[1][2.5cm]{%
  \nju@chapter*[]{\nju@denotation@name} % no tocline
  \vskip-30bp\xiaosi[1.6]\begin{nju@denotation}[labelwidth=#1]
}{%
  \end{nju@denotation}
}
\newlist{nju@denotation}{description}{1}
\setlist[nju@denotation]{%
  nosep,
  font=\normalfont,
  align=left,
  leftmargin=!, % sum of the following 3 lengths
  labelindent=0pt,
  labelwidth=2.5cm,
  labelsep*=0.5cm,
  itemindent=0pt,
}
%    \end{macrocode}
% \end{environment}
%
% \subsubsection{致谢与声明}
% \label{sec:ackanddeclare}
%
% \begin{environment}{acknowledgement}
% 支持扫描文件替换。
%    \begin{macrocode}
\ifnju@chinese
  \newcommand\nju@ack@name{致\hspace{\ccwd}谢}
\else
  \newcommand\nju@ack@name{Acknowledgments}
\fi
\newcommand\nju@declarename{声\hspace{\ccwd}明}
\newcommand{\nju@declaretext}{本人郑重声明：所呈交的学位论文，是本人在导师指导下
  ，独立进行研究工作所取得的成果。尽我所知，除文中已经注明引用的内容外，本学位论
  文的研究成果不包含任何他人享有著作权的内容。对本论文所涉及的研究工作做出贡献的
  其他个人和集体，均已在文中以明确方式标明。}
\newcommand{\nju@signature}{签\hspace{1em}名：}
\newcommand{\nju@backdate}{日\hspace{1em}期：}
%    \end{macrocode}
%
%  \cs{cleardoublepage}，使新开章节的页码到达正确的状态。否则会因为 \cs{addcontentsline}
% 在 chapter 之前而导致目录页码错误。
% 定义致谢与声明环境。
%    \begin{macrocode}
\NewDocumentEnvironment{acknowledgement}{o}{%
    \nju@chapter*{\nju@ack@name}
  }
%    \end{macrocode}
%
% 声明部分
%    \begin{macrocode}
  {
    \ifnju@english\relax\else%
      \IfNoValueTF{#1}{%
        \nju@chapter*{\nju@declarename}
        \par{\xiaosi\parindent2em\nju@declaretext}\vskip2cm
        {\xiaosi\hfill\nju@signature\nju@underline[2.5cm]\relax%
         \nju@backdate\nju@underline[2.5cm]\relax}%
      }{%
        \includepdf[pagecommand={\thispagestyle{nju@empty}%
          \phantomsection\addcontentsline{toc}{chapter}{\nju@declarename}%
        }]{#1}%
      }%
    \fi
  }
%    \end{macrocode}
% \end{environment}
%
% \subsubsection{图表索引}
% \label{sec:threeindex}
% \begin{macro}{\listoffigures}
% \begin{macro}{\listoffigures*}
% \begin{macro}{\listoftables}
% \begin{macro}{\listoftables*}
% 定义图表以及公式目录样式。
%    \begin{macrocode}
\def\nju@starttoc#1{% #1: float type, prepend type name in \listof*** entry.
  \let\oldnumberline\numberline
  \def\numberline##1{\oldnumberline{\csname #1name\endcsname\hskip.4em ##1}}
  \@starttoc{\csname ext@#1\endcsname}
  \let\numberline\oldnumberline}
\def\nju@listof#1{% #1: float type
  \@ifstar
    {\nju@chapter*[]{\csname list#1name\endcsname}\nju@starttoc{#1}}
    {\nju@chapter*{\csname list#1name\endcsname}\nju@starttoc{#1}}}
\renewcommand\listoffigures{\nju@listof{figure}}
\renewcommand*\l@figure{\addvspace{6bp}\@dottedtocline{1}{0em}{4em}}
\renewcommand\listoftables{\nju@listof{table}}
\let\l@table\l@figure
%    \end{macrocode}
% \end{macro}
% \end{macro}
% \end{macro}
% \end{macro}
%
% \begin{macro}{\equcaption}
%   本命令只是为了生成公式列表，所以这个 caption 是假的。如果要编号最好用
%    equation 环境，如果是其它编号环境，请手动添加 \cs{equcaption}。
% 用法如下：
%
% \cs{equcaption}\marg{counter}
%
% \marg{counter} 指定出现在索引中的编号，一般取 \cs{theequation}，如果你是用
%  \pkg{amsmath} 的 \cs{tag}，那么默认是 \cs{tag} 的参数；除此之外可能需要你
% 手工指定。
%
%    \begin{macrocode}
\def\ext@equation{loe}
\def\equcaption#1{%
  \addcontentsline{\ext@equation}{equation}%
                  {\protect\numberline{#1}}}
%    \end{macrocode}
% \end{macro}
%
% \begin{macro}{\listofequations}
% \begin{macro}{\listofequations*}
% \LaTeX\ 默认没有公式索引，此处定义自己的 \cs{listofequations}。
%    \begin{macrocode}
\newcommand\listofequations{\nju@listof{equation}}
\let\l@equation\l@figure
%    \end{macrocode}
% \end{macro}
% \end{macro}
%
% \subsection{参考文献}
% \label{sec:ref}
%
% \begin{macro}{\inlinecite}
% 依赖于 \pkg{natbib} 宏包，修改其中的命令。 旧命令 \cs{onlinecite} 依然可用。
%    \begin{macrocode}
\newcommand\bibstyle@inline{\bibpunct{[}{]}{,}{n}{,}{,}}
\DeclareRobustCommand\inlinecite{\@inlinecite}
\def\@inlinecite#1{\begingroup\let\@cite\NAT@citenum\citep{#1}\endgroup}
\let\onlinecite\inlinecite
%    \end{macrocode}
% \end{macro}
%
% 参考文献的正文部分用五号字。
% 行距采用固定值 16 磅，段前空 3 磅，段后空 0 磅。
%
% 复用 \pkg{natbib} 的 \texttt{thebibliography} 环境，调整距离。
%    \begin{macrocode}
\renewcommand\bibsection{\nju@chapter*{\bibname}}
\renewcommand\bibfont{\wuhao[1.5]}
\setlength\bibhang{2\ccwd}
\addtolength{\bibsep}{-0.7em}
\setlength{\labelsep}{0.4em}
\def\@biblabel#1{[#1]\hfill}
%    \end{macrocode}
%
% 两种引用样式：
%    \begin{macrocode}
\expandafter\newcommand\csname bibstyle@numeric\endcsname{%
  \bibpunct{[}{]}{,}{s}{,}{\textsuperscript{,}}}
\expandafter\newcommand\csname bibstyle@author-year\endcsname{%
  \bibpunct{(}{)}{;}{a}{,}{,}}
%    \end{macrocode}
%
% 下面修改 \pkg{natbib} 的引用格式，主要是将页码写在上标位置。
% numeric 模式的 \cs{citet} 的页码：
%    \begin{macrocode}
\patchcmd\NAT@citexnum{%
  \@ifnum{\NAT@ctype=\z@}{%
    \if*#2*\else\NAT@cmt#2\fi
  }{}%
  \NAT@mbox{\NAT@@close}%
}{%
  \NAT@mbox{\NAT@@close}%
  \@ifnum{\NAT@ctype=\z@}{%
    \if*#2*\else\textsuperscript{#2}\fi
  }{}%
}{}{}
%    \end{macrocode}
%
% Numeric 模式的 \cs{citep} 的页码：
%    \begin{macrocode}
\renewcommand\NAT@citesuper[3]{\ifNAT@swa
  \if*#2*\else#2\NAT@spacechar\fi
\unskip\kern\p@\textsuperscript{\NAT@@open#1\NAT@@close\if*#3*\else#3\fi}%
   \else #1\fi\endgroup}
%    \end{macrocode}
%
% Author-year 模式的 \cs{citet} 的页码：
%    \begin{macrocode}
\patchcmd{\NAT@citex}{%
  \if*#2*\else\NAT@cmt#2\fi
  \if\relax\NAT@date\relax\else\NAT@@close\fi
}{%
  \if\relax\NAT@date\relax\else\NAT@@close\fi
  \if*#2*\else\textsuperscript{#2}\fi
}{}{}
%    \end{macrocode}
%
% Author-year 模式的 \cs{citep} 的页码：
%    \begin{macrocode}
\renewcommand\NAT@citesuper[3]{\ifNAT@swa
  \if*#2*\else#2\NAT@spacechar\fi
\unskip\kern\p@\textsuperscript{\NAT@@open#1\NAT@@close\if*#3*\else#3\fi}%
   \else #1\fi\endgroup}
%    \end{macrocode}
%
% 在顺序编码制下，\pkg{natbib} 只有在三个以上连续文献引用才会使用连接号，
% 这里修改为允许两个引用使用连接号。
%    \begin{macrocode}
\patchcmd{\NAT@citexnum}{%
  \ifx\NAT@last@yr\relax
    \def@NAT@last@yr{\@citea}%
  \else
    \def@NAT@last@yr{--\NAT@penalty}%
  \fi
}{%
  \def@NAT@last@yr{-\NAT@penalty}%
}{}{}
%    \end{macrocode}
%
% \subsection{附录}
% \label{sec:appendix}
% \begin{environment}{appendix}
% 主要给本科做外文翻译用。
%    \begin{macrocode}
\let\nju@appendix\appendix
\renewenvironment{appendix}{%
  \let\title\nju@appendix@title
  \nju@appendix}{%
  \let\title\@gobble}
%    \end{macrocode}
% \end{environment}
%
% \begin{macro}{\title}
% 本科外文翻译文章的标题，用法：\cs{title}\marg{资料标题}。这个命令只能在附录环
% 境下使用。
%    \begin{macrocode}
\let\title\@gobble
\newcommand{\nju@appendix@title}[1]{%
  \begin{center}
    \xiaosi[1.667] #1
  \end{center}}
%    \end{macrocode}
% \end{macro}
%
% \begin{environment}{translationbib}
% 外文资料的参考文用宋体五号字，取固定行距17pt，段前后3pt。
%    \begin{macrocode}
\newlist{translationbib}{enumerate}{1}
\setlist[translationbib]{label=[\arabic*],align=left,nosep,itemsep=6bp,
  leftmargin=10mm,labelsep=!,before=\vspace{0.5\baselineskip}\wuhao[1.3]}
%    \end{macrocode}
% \end{environment}
%\marginpar{这是边注}
%\subsection{代码}
%    \begin{macrocode}
\RequirePackage{xcolor}
\definecolor{codegreen}{rgb}{0,0.6,0}
\definecolor{codegray}{rgb}{0.5,0.5,0.5}
\definecolor{codepurple}{rgb}{0.58,0,0.82}
\definecolor{backcolour}{rgb}{0.95,0.95,0.92}
\newcommand{\red}[1]{\textcolor{red}{#1}}
\newcommand{\redoverlay}[2]{\textcolor<#2>{red}{#1}}
\newcommand{\green}[1]{\textcolor{green}{#1}}
\newcommand{\greenoverlay}[2]{\textcolor<#2>{green}{#1}}
\newcommand{\blue}[1]{\textcolor{blue}{#1}}
\newcommand{\blueoverlay}[2]{\textcolor<#2>{blue}{#1}}
\newcommand{\purple}[1]{\textcolor{purple}{#1}}
\newcommand{\cyan}[1]{\textcolor{cyan}{#1}}
\newcommand{\teal}[1]{\textcolor{teal}{#1}}
\RequirePackage{verbatim}
\RequirePackage{algorithm}
\RequirePackage{algpseudocode}
\newcommand{\pseduo}[2]{
\begin{algorithm}
	\caption{\textsc{#1}}
	\label{alg:#1}
	\begin{algorithmic}[1]
		#2
	\end{algorithmic}
\end{algorithm}
}
\RequirePackage{listings}
\lstdefinestyle{lstStyleBase}{%
   basicstyle=\small\ttfamily,
   aboveskip=\medskipamount,
   belowskip=\medskipamount,
   lineskip=0pt,
   boxpos=c,
   showlines=false,
   extendedchars=true,
   upquote=true,
   tabsize=2,
   showtabs=false,
   showspaces=false,
   showstringspaces=false,
   numbers=none,
   linewidth=\linewidth,
   xleftmargin=4pt,
   xrightmargin=0pt,
   resetmargins=false,
   breaklines=true,
   breakatwhitespace=false,
   breakindent=0pt,
   breakautoindent=true,
   columns=flexible,
   keepspaces=true,
   gobble=2,
   framesep=3pt,
   rulesep=1pt,
   framerule=1pt,
   backgroundcolor=\color{gray!5},
   stringstyle=\color{green!40!black!100},
   keywordstyle=\bfseries\color{blue!50!black},
   commentstyle=\slshape\color{black!60}
}

\newtcblisting{commandshell}{colback=black,colupper=white,colframe=yellow!75!black, listing only,listing options={style=tcblatex,language=sh},
every listing line={\textcolor{red}{\small\ttfamily\bfseries root \$> }}}

\lstdefinestyle{lstStyleShell}{%
   style=lstStyleBase,
   frame=l,
   rulecolor=\color{purple},
   language=bash}

\lstdefinestyle{lstStyleLaTeX}{%
   style=lstStyleBase,
   frame=l,
   rulecolor=\color{violet},
   language=[LaTeX]TeX}

\lstdefinestyle{lstStylecdisplay}{%
  style=lstStyleBase,
  frame=tb,
  rulecolor=\color{cyan},
  keywordstyle=\color{magenta}\bfseries\ttfamily,
  commentstyle=\color{codegreen}\ttfamily,
	stringstyle=\color{codepurple}\ttfamily\sffamily,
	backgroundcolor=\color{backcolour},
	captionpos=b,
	numbers=left,
	numberstyle=\footnotesize\color{codegray},
	stepnumber=1,
  numbersep=5pt,
  language=C
}

\lstdefinestyle{lstStylecpseudo}{%
  style=lstStyleBase,
  frame=none,
  keywordstyle=\color{magenta}\bfseries\ttfamily,
  commentstyle=\color{codegreen}\ttfamily,
	stringstyle=\color{codepurple}\ttfamily\sffamily,
	captionpos=b,
	numbers=left,
	numberstyle=\footnotesize\color{codegray},
	stepnumber=1,
  numbersep=5pt,
  language=C
}

\lstdefinestyle{lstStylecplus}{%
  style=lstStyleBase,
  frame=l,
  rulecolor=\color{blue},
  language=C++
}

\lstdefinestyle{lstStyleverilog}{%
  style=lstStyleBase,
  frame=l,
  rulecolor=\color{brown},
  language=verilog
}

\lstdefinestyle{lstStylepython}{%
  style=lstStyleBase,
  frame=l,
  rulecolor=\color{pink},
  language=python
}

\lstnewenvironment{code}{\lstset{style=lstStyleBase}}{}
\lstnewenvironment{latex}{\lstset{style=lstStyleLaTeX}}{}
\lstnewenvironment{shell}{\lstset{style=lstStyleShell}}{}
\lstnewenvironment{cdisplay}{\lstset{style=lstStylecdisplay}}{}
\lstnewenvironment{cplus}{\lstset{style=lstStylecplus}}{}
\lstnewenvironment{verilog}{\lstset{style=lstStyleverilog}}{}
\lstnewenvironment{python}{\lstset{style=lstStylepython}}{}
\lstnewenvironment{cpseudo}{\lstset{style=lstStylecpseudo}}{}
%    \end{macrocode}
%
% \subsection{快速插入图片或图表}
%    \begin{macrocode}
\newcommand{\figpf}[2]{
	\begin{figure}[H]
		\centering
		\includegraphics[#1]{figs/#2}
	\end{figure}
}

%%%%%%%%%%%%%%%%%%%%
\newcommand{\figpfc}[3]{
	\begin{figure}[htbp]
		\centering
		\includegraphics[#1]{figs/#2}
		\caption{#3}
		\label{fig:#2}
	\end{figure}
}
%%%%%%%%%%%%%%%%%%%
\newcommand{\tabpcc}[3]{
	\begin{table}[H]
		\centering
		\begin{tabular}{|*{#1}{c|}}
		\toprule
		#2\\
		\bottomrule
	\end{tabular}
	\caption{#3}
	\label{form:#3}
\end{table}}
%%%%%%%%%%%%%%%%%%%
\newcommand{\tabpc}[2]{
	\begin{table}[H]
		\centering
		\begin{tabular}{|*{#1}{c|}}
		\toprule
		#2\\
		\bottomrule
	\end{tabular}
\end{table}}
\newcommand{\tnl}{\tabularnewline\midrule}
%    \end{macrocode}
% \subsection{结束部分}
% \label{sec:finish}
%    \begin{macrocode}
\def\cmd#1{\cs{\expandafter\cmd@to@cs\string#1}}
\def\cmd@to@cs#1#2{\char\number`#2\relax}
\DeclareRobustCommand\cs[1]{\texttt{\char`\\#1}}

\newcommand*{\meta}[1]{{%
  \ensuremath{\langle}\rmfamily\itshape#1\/\ensuremath{\rangle}}}
\providecommand\marg[1]{%
  {\ttfamily\char`\{}\meta{#1}{\ttfamily\char`\}}}
\providecommand\oarg[1]{%
  {\ttfamily[}\meta{#1}{\ttfamily]}}
\providecommand\parg[1]{%
  {\ttfamily(}\meta{#1}{\ttfamily)}}
\providecommand\pkg[1]{{\sffamily#1}}


% 表格中支持跨行
\RequirePackage{multirow}
\RequirePackage{bbding,stmaryrd}
% 固定宽度的表格
\RequirePackage{tabularx}
\RequirePackage{syntonly}
\RequirePackage{dirtree}
% 表格中的反斜线
\RequirePackage{diagbox}
\newcommand{\blankpage}{
	\clearpage
	\begin{titlepage}
		\null\vfil
		\begin{center}
			\textit{This page intentionally left blank.}
		\end{center}
	\end{titlepage}
}\RequirePackage{syntonly}
\newcommand{\rmnum}[1]{\romannumeral #1}
\newcommand{\Rmnum}[1]{\expandafter\@slowromancap\romannumeral #1@}
% 确定浮动对象的位置，可以使用 H，强制将浮动对象放到这里（可能效果很差）
\RequirePackage{makecell}
\RequirePackage{physics}
\RequirePackage{float}
\RequirePackage{watermark} 
\RequirePackage{tikz,tikzducks}
\usetikzlibrary{decorations.pathmorphing,graphs,calc}
\ifnju@draft
\AtEndOfClass{
	\watermark{% 
		\parbox[b][\paperheight]{\paperwidth}{% 
		\vfill 
		\centering% 
		\begin{tikzpicture}[remember picture,overlay] 
			\node [rotate=45,scale=10] at ($(current page.center) +(-1cm,1cm)$) 
			{\textcolor[gray]{0.8}{DRAFT}}; 
			\node [rotate=45,scale=3] at ($(current page.center) +(1cm,-1cm)$) 
			{\textcolor[gray]{0.75}{Compile time: \the\year - \the\month - \the\day}}; 
		\end{tikzpicture}% 
		\vfill 
		}
  }
}
\fi

\AtEndOfClass{\sloppy}
%    \end{macrocode}
%</cls> 
%
%
%
%
%
%
%
% \iffalse
%    \begin{macrocode}
%<*dtx-style>
\ProvidesPackage{dtx-style}
\RequirePackage{hypdoc}
\RequirePackage{ifthen}
\RequirePackage[UTF8,scheme=chinese]{ctex}
\RequirePackage{newpxtext}
\RequirePackage{newpxmath}
\RequirePackage[
  top=2.5cm, bottom=2.5cm,
  left=4cm, right=2cm,marginparwidth=2.6cm,marginparsep=3mm,
  headsep=3mm]{geometry}
\RequirePackage{array,longtable,booktabs}
\RequirePackage{listings}
\RequirePackage{fancyhdr}
\RequirePackage{xcolor}
\definecolor{codegreen}{rgb}{0,0.6,0}
\definecolor{codegray}{rgb}{0.5,0.5,0.5}
\definecolor{codepurple}{rgb}{0.58,0,0.82}
\definecolor{backcolour}{rgb}{0.95,0.95,0.92}
\newcommand{\red}[1]{\textcolor{red}{#1}}
\newcommand{\redoverlay}[2]{\textcolor<#2>{red}{#1}}
\newcommand{\green}[1]{\textcolor{green}{#1}}
\newcommand{\greenoverlay}[2]{\textcolor<#2>{green}{#1}}
\newcommand{\blue}[1]{\textcolor{blue}{#1}}
\newcommand{\blueoverlay}[2]{\textcolor<#2>{blue}{#1}}
\newcommand{\purple}[1]{\textcolor{purple}{#1}}
\newcommand{\cyan}[1]{\textcolor{cyan}{#1}}
\newcommand{\teal}[1]{\textcolor{teal}{#1}}
\RequirePackage{enumitem}
\RequirePackage{etoolbox}
\RequirePackage{metalogo}
\RequirePackage{mathtools}
\DeclarePairedDelimiter{\ceil}{\lceil}{\rceil}
\DeclarePairedDelimiter{\floor}{\lfloor}{\rfloor}
\DeclareMathOperator{\Hamilton}{\hat{H}} 
\ifthenelse{\equal{\@nameuse{g__ctex_fontset_tl}}{mac}}{%
  \xeCJKsetwidth{‘’“”}{1em}
}{}

\colorlet{nju@macro}{blue!60!black}
\colorlet{nju@env}{blue!70!black}
\colorlet{nju@option}{purple}
\patchcmd{\PrintMacroName}{\MacroFont}{\MacroFont\bfseries\color{nju@macro}}{}{}
\patchcmd{\PrintDescribeMacro}{\MacroFont}{\MacroFont\bfseries\color{nju@macro}}{}{}
\patchcmd{\PrintDescribeEnv}{\MacroFont}{\MacroFont\bfseries\color{nju@env}}{}{}
\patchcmd{\PrintEnvName}{\MacroFont}{\MacroFont\bfseries\color{nju@env}}{}{}

\def\DescribeOption{%
  \leavevmode\@bsphack\begingroup\MakePrivateLetters%
  \Describe@Option}
\def\Describe@Option#1{\endgroup
  \marginpar{\raggedleft\PrintDescribeOption{#1}}%
  \nju@special@index{option}{#1}\@esphack\ignorespaces}
\def\PrintDescribeOption#1{\strut \MacroFont\bfseries\sffamily\color{nju@option} #1\ }
\def\nju@special@index#1#2{\@bsphack
  \begingroup
    \HD@target
    \let\HDorg@encapchar\encapchar
    \edef\encapchar usage{%
      \HDorg@encapchar hdclindex{\the\c@HD@hypercount}{usage}%
    }%
    \index{#2\actualchar{\string\ttfamily\space#2}
           (#1)\encapchar usage}%
    \index{#1:\levelchar#2\actualchar
           {\string\ttfamily\space#2}\encapchar usage}%
  \endgroup
  \@esphack}

\lstdefinestyle{lstStyleBase}{%
   basicstyle=\small\ttfamily,
   aboveskip=\medskipamount,
   belowskip=\medskipamount,
   lineskip=0pt,
   boxpos=c,
   showlines=false,
   extendedchars=true,
   upquote=true,
   tabsize=2,
   showtabs=false,
   showspaces=false,
   showstringspaces=false,
   numbers=none,
   linewidth=\linewidth,
   xleftmargin=4pt,
   xrightmargin=0pt,
   resetmargins=false,
   breaklines=true,
   breakatwhitespace=false,
   breakindent=0pt,
   breakautoindent=true,
   columns=flexible,
   keepspaces=true,
   gobble=2,
   framesep=3pt,
   rulesep=1pt,
   framerule=1pt,
   backgroundcolor=\color{gray!5},
   stringstyle=\color{green!40!black!100},
   keywordstyle=\bfseries\color{blue!50!black},
   commentstyle=\slshape\color{black!60}}

\lstdefinestyle{lstStyleShell}{%
   style=lstStyleBase,
   frame=l,
   rulecolor=\color{purple},
   language=bash}

\lstdefinestyle{lstStyleLaTeX}{%
   style=lstStyleBase,
   frame=l,
   rulecolor=\color{violet},
   language=[LaTeX]TeX}

\lstnewenvironment{latex}{\lstset{style=lstStyleLaTeX}}{}
\lstnewenvironment{shell}{\lstset{style=lstStyleShell}}{}

\setlist{nosep}

\DeclareDocumentCommand{\option}{m}{\textsf{#1}}
\DeclareDocumentCommand{\env}{m}{\texttt{#1}}
\DeclareDocumentCommand{\pkg}{s m}{%
  \texttt{#2}\IfBooleanF#1{\nju@special@index{package}{#2}}}
\DeclareDocumentCommand{\file}{s m}{%
  \texttt{#2}\IfBooleanF#1{\nju@special@index{file}{#2}}}
\newcommand{\myentry}[1]{%
  \marginpar{\raggedleft\color{purple}\bfseries\strut #1}}
\newcommand{\note}[2][Note]{{%
  \color{magenta}{\bfseries #1}\emph{#2}}}

\def\njurepo{\textsc{NJU}\-\textsc{repo}}
\def\thuthesis{\textsc{Thu}\-\textsc{Thesis}}
%</dtx-style>
%    \end{macrocode}
% \fi
% \Finale
